\uuid{zfOt}
\titre{ Nettoyage de données : ordre des opérations et choix de la stratégie }
\niveau{L3}
\module{Analyse de données}
\chapitre{Prétraitement des données}
\sousChapitre{Nettoyage, valeurs manquantes, filtrage}
\theme{dropna, fillna, valeurs aberrantes, ordre des opérations}
\auteur{Maxime Nguyen}
\datecreate{2026-03-19}
\organisation{AMSCC}
\difficulte{2}

\contenu{
\texte{
  On compare trois façons de nettoyer un DataFrame \texttt{df} contenant des colonnes \texttt{SqFt} et \texttt{PRV} :
  \begin{itemize}
    \item[A.] \texttt{df = df.dropna()} puis \texttt{df = df[(df["SqFt"] > 0) \& (df["PRV"] > 0)]}
    \item[B.] \texttt{df = df[(df["SqFt"] > 0) \& (df["PRV"] > 0)]} puis \texttt{df = df.dropna()}
    \item[C.] \texttt{df = df.fillna(0)} puis \texttt{df = df[(df["SqFt"] > 0) \& (df["PRV"] > 0)]}
  \end{itemize}
}
\begin{enumerate}

\item
  \question{A et B produisent-ils toujours le même résultat ? Donner un exemple où ils diffèrent.}
  \reponse{
    Non. Considérons une ligne avec \texttt{SqFt = NaN} et \texttt{PRV = 500}.
    \begin{itemize}
      \item \textbf{A} supprime cette ligne dès le \texttt{dropna()}, donc elle n'arrive pas au filtre.
      \item \textbf{B} évalue d'abord \texttt{df["SqFt"] > 0} : \texttt{NaN > 0} vaut \texttt{False},
            donc la ligne est supprimée par le filtre, puis \texttt{dropna()} ne trouve plus rien à faire.
    \end{itemize}
    Le résultat final est identique dans ce cas, mais diffère si une ligne a \texttt{SqFt = NaN}
    et \texttt{PRV = NaN} : avec B, la comparaison \texttt{NaN > 0} donne \texttt{False},
    la ligne est retirée par le filtre ; avec A, \texttt{dropna()} la retire en premier.
    En pratique, A et B donnent souvent le même résultat, mais A est plus lisible et explicite.
  }

\item
  \question{Quel problème pose C par rapport à A ?}
  \reponse{
    C remplace les valeurs manquantes par $0$ avant de filtrer.
    Or \texttt{SqFt = 0} et \texttt{PRV = 0} ne satisfont pas \texttt{> 0},
    donc ces lignes sont ensuite supprimées — ce qui revient à supprimer les valeurs manquantes
    mais en introduisant des $0$ parasites pour les autres colonnes non filtrées.
    De plus, imputer $0$ pour une surface (\texttt{SqFt}) ou un prix (\texttt{PRV}) est sémantiquement faux
    et peut biaiser les modèles si d'autres colonnes avec \texttt{NaN} reçoivent aussi $0$.
  }

\end{enumerate}
}
